% Inactive research fragment; no general bounded-to-smooth theorem is asserted.
% Inputs: interior Lp Laplace estimates, Sobolev compactness, and the declared
% Garding algebra in the preceding fragments. No full real Hessian bound is used.
\section{Limits with bounded complex Hessians}

Let $(X,\omega)$ be a compact Hermitian manifold of dimension $n$.
In a coordinate chart write
\[
 \omega=\frac{\mathrm i}{2\pi}g_{i\bar j}\,
                 \mathrm dz_i\wedge\mathrm d\bar z_j,
 \qquad H(u)=(u_{i\bar j}).
\]
The norms of complex Hessians below are measured with respect to the fixed
metric. The pointwise coefficient products of forms with bounded measurable
coefficients will be used only after their distributional identities have
been justified by local convolution.

\begin{lemma}\label{draft:603-complex-hessian-compactness}
Suppose $u_j$ is smooth,
\[
 \sup_j\bigl(\|u_j\|_\infty+\|\ddc u_j\|_{L^\infty}\bigr)<\infty,
 \qquad u_j\longrightarrow u\quad\hbox{uniformly on }X.
\]
Then $u\in W^{2,p}(X)$ for every $1<p<\infty$, its distributional complex
Hessian has bounded coefficients, and $u_j\to u$ in $C^{1,\alpha}(X)$
for every $0<\alpha<1$. The same hypotheses give uniform $W^{2,p}$ bounds
for the sequence.
\end{lemma}
\begin{proof}
On a coordinate ball the Euclidean Laplacian is
$\Delta u_j=4\sum_i (u_j)_{i\bar i}$, so it is uniformly bounded.
The interior Laplace estimate on two nested balls gives
\[
 \|u_j\|_{W^{2,p}(B')}\leq C_p
       \bigl(\|\Delta u_j\|_{L^p(B)}+\|u_j\|_{L^p(B)}\bigr).
\]
A finite cover gives the global bound. Weak $W^{2,p}$ limits coincide
with the uniform limit. The bounded complex Hessians have weak-star
limits in $L^\infty$; testing against compactly supported smooth functions
identifies them with $H(u)$. Choose $p$ with $1-2n/p>\alpha$.
The compact Sobolev embedding into $C^{1,\alpha}$ gives convergence of
subsequences. Every such limit is $u$, hence the whole sequence converges.
This estimates the full real Hessian in every finite $L^p$ space, without
asserting that it is bounded.
\end{proof}

\begin{lemma}\label{draft:603-complex-hessian-products}
For $1\leq a\leq k\leq n$, suppose $u_j^a$ and $u^a$ satisfy the preceding
lemma. Let $\theta_a$ be fixed smooth real $(1,1)$-forms. Then
\[
 \bigwedge_{a=1}^k(\theta_a+\ddc u_j^a)
       \longrightarrow
 \bigwedge_{a=1}^k(\theta_a+\ddc u^a)
       \quad\hbox{as currents}.
\]
The current on the right is the product of the almost-everywhere
coefficients. No closedness or positivity of the $\theta_a$ is required.
For each compactly supported smooth complementary test form $\eta$, the
difference of its pairings is bounded by
$C_\eta\max_a\|u_j^a-u^a\|_\infty$.
\end{lemma}
\begin{proof}
Set $D_j^a=\ddc u_j^a$ and $D^a=\ddc u^a$.
On a smaller coordinate ball, convolution gives smooth $u_\epsilon^a$
whose complex Hessians are uniformly bounded and converge to $H(u^a)$
in every finite $L^p$ space. Their coefficient products consequently
converge in finite $L^p$ spaces. In particular the products of the $D^a$
are closed currents, since their smooth approximants are closed.

Expand the smooth background factors and telescope each remaining
product. Every difference has the form
\[
 \ddc(u_j^a-u^a)\wedge T_j,
\]
where $T_j$ is a product of $D_j^b$ and $D^b$ with $b\ne a$.
It has uniformly bounded coefficients and is closed.
Local convolution, followed by strong convergence of its finite
coefficient products, justifies integration by parts:
\[
 \left\langle\ddc(u_j^a-u^a)\wedge T_j,\xi\right\rangle
       =\int (u_j^a-u^a)\,T_j\wedge\ddc\xi.
\]
Here $\xi$ includes the chosen smooth background factors and the test
form, so their derivatives are retained. The right-hand side has the
asserted uniform bound. A finite partition of unity proves the claim on
$X$. This argument also proves that the coefficient product is the
distributional limit; it does not identify a separately defined bounded
Hessian operator for all nonsmooth admissible functions.
\end{proof}

Let $2\leq m\leq n$, let $C_{n,m}$ be the closed Hermitian G\aa rding cone,
and put $F(A)=\sigma_m(\lambda(A))^{1/m}$ in its interior, with
$\sigma_m=S_m/\binom nm$. We use its classical concavity and Maclaurin
inequalities, and the derivative estimate proved in
\cref{draft:599-weighted-cone-trace}.

\begin{lemma}\label{draft:603-uniform-elliptic-extension}
For $f_0>0$ and $L>0$, assume the set
\[
 K=\{A\in C_{n,m}:\operatorname{tr}A\leq L,
                         \ \sigma_m(A)\geq f_0^m\}
\]
is nonempty. It is a compact subset of the interior, and there are
$\vartheta,\Theta>0$, depending only on $n,m,f_0,L$, such that
\[
 \vartheta I\leq DF_A\leq\Theta I\qquad(A\in K).
\]
The function on all Hermitian matrices defined by
\[
 \widehat F(B)=\inf_{A\in K}
       \bigl(F(A)+\operatorname{tr}(DF_A(B-A))\bigr)
\]
is concave, agrees with $F$ on $K$, and satisfies, for $P\geq0$,
\[
 \vartheta\operatorname{tr}P\leq
       \widehat F(B+P)-\widehat F(B)
       \leq\Theta\operatorname{tr}P.
\]
\end{lemma}
\begin{proof}
In the closed $2$-cone,
\[
 \sum_i\lambda_i^2=(\operatorname{tr}A)^2-2S_2(A)
                          \leq L^2.
\]
Thus $K$ is bounded and closed. Maclaurin gives
$\sigma_j(A)\geq\sigma_m(A)^{j/m}\geq f_0^j$ for $j\leq m$,
so every matrix in $K$ is in the interior. In a unitary eigenframe,
\[
 F_i=\frac{F^{1-m}S_{m-1}(\lambda|i)}{m\binom nm},
 \qquad
 \frac{f_0}{m!L}\leq F_i\leq
                 \frac{L^{m-1}}{n f_0^{m-1}}.
\]
The lower bound is the preceding derivative estimate; the upper bound
uses $|\lambda_i|\leq L$ in the polynomial numerator and its positivity.
These bounds give possible values of $\vartheta$ and $\Theta$.
Concavity places each tangent plane above $F$ on the cone; the plane
at $B$ gives equality when $B\in K$. An infimum of affine functions is
concave. Each plane increases under $P\geq0$ by a number between the
two displayed trace bounds; taking infima preserves these bounds.
The same estimates give global Lipschitz continuity of $\widehat F$.
\end{proof}

To interpret this as a real elliptic operator, let $\mathcal C(R)$ be
the complex Hessian of the real quadratic function
$q(x)=\frac12 x^{\mathsf T}Rx$, where $x\in\mathbb R^{2n}$.
For $R\geq0$, one has $\mathcal C(R)\geq0$ and
$\operatorname{tr}\mathcal C(R)=\operatorname{tr}R/4$.
If $\mu I\leq g(x)\leq MI$, then
\[
 \mathcal G(x,R)=\widehat F
       \bigl(I+g(x)^{-1/2}\mathcal C(R)g(x)^{-1/2}\bigr)
\]
is uniformly elliptic on real symmetric matrices, with constants
$\vartheta/(4M)$ and $\Theta/(4\mu)$.
Its dependence on $x$ satisfies on compact coordinate balls
\[
 |\mathcal G(x,R)-\mathcal G(y,R)|
                      \leq C|x-y|(1+\|R\|).
\]
These follow respectively from the trace bound and the smooth dependence
of the positive matrix square root of the fixed metric.

\begin{proposition}\label{draft:603-proper-auxiliary-limit}
Suppose $v_j\to\varphi$ uniformly on $X$, $h(x,t)$ is smooth on a fixed
bounded interval containing their values, and smooth admissible $u_j$
satisfy
\[
 (\omega+\ddc u_j)^m\wedge\omega^{n-m}
           =e^{u_j+h(x,v_j)}\omega^n.
\]
Assume $u_j\to u$ uniformly and the complex Hessians of $u_j$ are uniformly
bounded. Then $\omega+\ddc u$ is admissible almost everywhere and its
coefficient product satisfies
\[
 (\omega+\ddc u)^m\wedge\omega^{n-m}
           =e^{u+h(x,\varphi)}\omega^n.
\]
On every coordinate ball $u$ also solves the corresponding uniformly
elliptic concave extension in the viscosity sense.
\end{proposition}
\begin{proof}
The first lemma gives $u\in W^{2,p}$ for every finite $p$ and a bounded
complex Hessian. The matrices
$A_j=I+g^{-1/2}H(u_j)g^{-1/2}$ converge weakly in $L^2$ to the analogous
matrix $A$ for $u$. Mazur's lemma gives convex combinations of tails
converging strongly in $L^2$, hence almost everywhere along a subsequence.
The pointwise closed cone is convex, so $A\in C_{n,m}$ almost everywhere.
The product lemma passes the equation to the limit as currents, since its
right-hand side converges uniformly. Equality of the bounded densities
then gives the equation almost everywhere.

The right-hand sides are bounded above and away from zero. The complex
Hessian bound gives a uniform upper bound $L$ for $\operatorname{tr}A_j$.
Choose $f_0>0$ below all their normalized roots. Thus every $A_j$, and
almost every $A$, belongs to the compact set $K$ of the preceding lemma.
Every $u_j$ is a classical solution of
\[
 \mathcal G(x,D^2u_j)=r_j(x),\qquad
               r_j(x)=e^{(u_j+h(x,v_j))/m}.
\]
For completeness, if a smooth test $q$ touches $u$ from above at an
interior point $x_0$, add a small positive quartic term to make the local
contact strict. Local maximum points of $u_j-q$ tend to $x_0$.
At these points $D^2u_j\leq D^2q$, and ellipticity gives
$\mathcal G(x,D^2q)\geq r_j(x)$. Pass to the limit, then remove the quartic
term. A test from below gives the reverse inequality. The extension is
defined on every real Hessian, so no admissibility of a lower test has
been assumed. This proves the viscosity statement.
\end{proof}

In the regularity application $h(x,t)=-t+f(x,t)$, identification of this
limit with $\varphi$ still requires bounded Hessian products for the given
weak solution and exponential uniqueness. Continuity of that weak
solution, existence of the smooth auxiliary solutions, global admissible
approximation, and the subsequent Evans--Krylov and Schauder applications
are separate dependencies. This fragment neither proves nor assumes that
an arbitrary bounded admissible weak solution is already continuous.
